% Options for packages loaded elsewhere
\PassOptionsToPackage{unicode}{hyperref}
\PassOptionsToPackage{hyphens}{url}
\documentclass[
  11pt,
  a4paper,
]{article}
\usepackage{xcolor}
\usepackage[margin=18mm]{geometry}
\usepackage{amsmath,amssymb}
\setcounter{secnumdepth}{-\maxdimen} % remove section numbering
\usepackage{iftex}
\ifPDFTeX
  \usepackage[T1]{fontenc}
  \usepackage[utf8]{inputenc}
  \usepackage{textcomp} % provide euro and other symbols
\else % if luatex or xetex
  \usepackage{unicode-math} % this also loads fontspec
  \defaultfontfeatures{Scale=MatchLowercase}
  \defaultfontfeatures[\rmfamily]{Ligatures=TeX,Scale=1}
\fi
\ifPDFTeX\else
  % xetex/luatex font selection
    \setmainfont[]{Noto Serif}
    \setsansfont[]{Noto Sans}
    \setmonofont[]{Noto Sans Mono}
  \setmathfont[]{Noto Sans Math}
\fi
% Use upquote if available, for straight quotes in verbatim environments
\IfFileExists{upquote.sty}{\usepackage{upquote}}{}
\IfFileExists{microtype.sty}{% use microtype if available
  \usepackage[]{microtype}
  \UseMicrotypeSet[protrusion]{basicmath} % disable protrusion for tt fonts
}{}
\makeatletter
\@ifundefined{KOMAClassName}{% if non-KOMA class
  \IfFileExists{parskip.sty}{%
    \usepackage{parskip}
  }{% else
    \setlength{\parindent}{0pt}
    \setlength{\parskip}{6pt plus 2pt minus 1pt}}
}{% if KOMA class
  \KOMAoptions{parskip=half}}
\makeatother
\ifLuaTeX
\usepackage[bidi=basic]{babel}
\else
\usepackage[bidi=default]{babel}
\fi
\babelprovide[main,import]{russian}
\ifPDFTeX
\else
\babelfont{rm}[]{Noto Serif}
\fi
% get rid of language-specific shorthands (see #6817):
\let\LanguageShortHands\languageshorthands
\def\languageshorthands#1{}
\setlength{\emergencystretch}{3em} % prevent overfull lines
\providecommand{\tightlist}{%
  \setlength{\itemsep}{0pt}\setlength{\parskip}{0pt}}
\usepackage{bookmark}
\IfFileExists{xurl.sty}{\usepackage{xurl}}{} % add URL line breaks if available
\urlstyle{same}
\hypersetup{
  pdftitle={Билеты по дискретным случайным процессам},
  pdflang={ru-RU},
  hidelinks,
  pdfcreator={LaTeX via pandoc}}

\title{Билеты по дискретным случайным процессам}
\author{}
\date{}

\begin{document}
\maketitle

\section{Билет 16. Дискретный случайный процесс как последовательность
случайных величин и как случайная
последовательность}\label{ux431ux438ux43bux435ux442-16.-ux434ux438ux441ux43aux440ux435ux442ux43dux44bux439-ux441ux43bux443ux447ux430ux439ux43dux44bux439-ux43fux440ux43eux446ux435ux441ux441-ux43aux430ux43a-ux43fux43eux441ux43bux435ux434ux43eux432ux430ux442ux435ux43bux44cux43dux43eux441ux442ux44c-ux441ux43bux443ux447ux430ux439ux43dux44bux445-ux432ux435ux43bux438ux447ux438ux43d-ux438-ux43aux430ux43a-ux441ux43bux443ux447ux430ux439ux43dux430ux44f-ux43fux43eux441ux43bux435ux434ux43eux432ux430ux442ux435ux43bux44cux43dux43eux441ux442ux44c}

Источники: Ширяев 2; Сердобольская, лекция 1.

\subsection{1. Короткий
ответ}\label{ux43aux43eux440ux43eux442ux43aux438ux439-ux43eux442ux432ux435ux442}

Дискретный случайный процесс можно понимать двумя эквивалентными
способами.

Первый способ: это семейство случайных величин

\[
\{X_n:n\in\mathbb N_0\}
\]

на одном вероятностном пространстве

\[
(\Omega,\mathcal F,P),
\]

где каждое

\[
X_n:\Omega\to E
\]

измеримо.

Второй способ: это одна случайная последовательность, то есть одно
измеримое отображение

\[
X:\Omega\to E^{\mathbb N_0},
\qquad
X(\omega)=(X_0(\omega),X_1(\omega),\ldots).
\]

Закон процесса - это распределение этого отображения:

\[
P_X(A)=P\{\omega:X(\omega)\in A\},
\qquad
A\in\mathcal E^{\mathbb N_0}.
\]

Конечномерные распределения процесса - это маргиналы закона \(P_X\) на
конечные наборы координат:

\[
P\{(X_{n_1},\ldots,X_{n_k})\in B\}
=
P_X\{x:(x_{n_1},\ldots,x_{n_k})\in B\}.
\]

\subsection{2. Полный
ответ}\label{ux43fux43eux43bux43dux44bux439-ux43eux442ux432ux435ux442}

\subsubsection{2.1. Введение и дискретное
время}\label{ux432ux432ux435ux434ux435ux43dux438ux435-ux438-ux434ux438ux441ux43aux440ux435ux442ux43dux43eux435-ux432ux440ux435ux43cux44f}

Этот билет продолжает Билет 15. Там теорема Колмогорова объясняла, как
по согласованным конечномерным распределениям построить меру на
пространстве траекторий. Здесь нужно разобрать самый важный для курса
частный случай: дискретное время, то есть случайные последовательности.

Пусть время дискретно:

\[
n=0,1,2,\ldots
\]

или

\[
n\in\mathbb Z.
\]

Обычно в этом билете достаточно рассматривать

\[
\mathbb N_0=\{0,1,2,\ldots\}.
\]

Пусть пространство состояний равно \(E\), например \(E=\mathbb R\),
\(E=\mathbb R^d\) или счетное множество состояний. На \(E\) задана
\(\sigma\)-алгебра \(\mathcal E\), обычно борелевская \(\mathcal B(E)\).

\subsubsection{2.2. Первая точка зрения: последовательность случайных
величин}\label{ux43fux435ux440ux432ux430ux44f-ux442ux43eux447ux43aux430-ux437ux440ux435ux43dux438ux44f-ux43fux43eux441ux43bux435ux434ux43eux432ux430ux442ux435ux43bux44cux43dux43eux441ux442ux44c-ux441ux43bux443ux447ux430ux439ux43dux44bux445-ux432ux435ux43bux438ux447ux438ux43d}

Первая точка зрения: дискретный случайный процесс - это
последовательность случайных элементов

\[
X_0,X_1,X_2,\ldots
\]

на одном вероятностном пространстве

\[
(\Omega,\mathcal F,P),
\]

где

\[
X_n:(\Omega,\mathcal F)\to(E,\mathcal E)
\]

измеримо для каждого \(n\).

Здесь \(n\) - время, а \(\omega\) - исход. Для фиксированного \(n\)
величина

\[
X_n(\omega)
\]

показывает состояние процесса в момент времени \(n\).

\subsubsection{2.3. Вторая точка зрения: пространство
траекторий}\label{ux432ux442ux43eux440ux430ux44f-ux442ux43eux447ux43aux430-ux437ux440ux435ux43dux438ux44f-ux43fux440ux43eux441ux442ux440ux430ux43dux441ux442ux432ux43e-ux442ux440ux430ux435ux43aux442ux43eux440ux438ux439}

Вторая точка зрения: весь процесс можно рассматривать как одно
отображение в пространство последовательностей:

\[
X:\Omega\to E^{\mathbb N_0},
\]

где

\[
X(\omega)=(X_0(\omega),X_1(\omega),X_2(\omega),\ldots).
\]

То есть каждому исходу \(\omega\) сопоставляется целая траектория
процесса.

Пространство последовательностей:

\[
E^{\mathbb N_0}
=
\{x=(x_0,x_1,x_2,\ldots):x_n\in E\}.
\]

На нем берут произведенную \(\sigma\)-алгебру

\[
\mathcal E^{\mathbb N_0}
=
\sigma\{\pi_n^{-1}(A):n\ge0,\ A\in\mathcal E\},
\]

где

\[
\pi_n:E^{\mathbb N_0}\to E,
\qquad
\pi_n(x)=x_n
\]

\begin{itemize}
\tightlist
\item
  координатная проекция.
\end{itemize}

\subsubsection{2.4. Конечномерные проекции и цилиндрические
множества}\label{ux43aux43eux43dux435ux447ux43dux43eux43cux435ux440ux43dux44bux435-ux43fux440ux43eux435ux43aux446ux438ux438-ux438-ux446ux438ux43bux438ux43dux434ux440ux438ux447ux435ux441ux43aux438ux435-ux43cux43dux43eux436ux435ux441ux442ux432ux430}

Более общие конечномерные проекции имеют вид

\[
\pi_{n_1,\ldots,n_k}:E^{\mathbb N_0}\to E^k,
\qquad
\pi_{n_1,\ldots,n_k}(x)=(x_{n_1},\ldots,x_{n_k}).
\]

Цилиндрическое множество - это событие, зависящее только от конечного
числа координат:

\[
C=\pi_{n_1,\ldots,n_k}^{-1}(B),
\qquad
B\in\mathcal E^k.
\]

Например,

\[
C=\{x:x_0\in A_0,\ x_1\in A_1,\ x_5\in A_5\}
\]

является цилиндром: он накладывает условия только на координаты
\(0,1,5\), а остальные координаты свободны.

\subsubsection{2.5. Доказательство эквивалентности точек
зрения}\label{ux434ux43eux43aux430ux437ux430ux442ux435ux43bux44cux441ux442ux432ux43e-ux44dux43aux432ux438ux432ux430ux43bux435ux43dux442ux43dux43eux441ux442ux438-ux442ux43eux447ux435ux43a-ux437ux440ux435ux43dux438ux44f}

Теперь важно показать эквивалентность двух точек зрения.

Если задана последовательность случайных величин \(X_n\), то определяем

\[
X(\omega)=(X_0(\omega),X_1(\omega),\ldots).
\]

Нужно проверить, что \(X\) измеримо как отображение в
\(E^{\mathbb N_0}\). Поскольку произведенная \(\sigma\)-алгебра
порождена цилиндрами, достаточно проверить прообразы цилиндров. Для
одного координатного цилиндра:

\[
X^{-1}(\pi_n^{-1}(A))
=
\{\omega:X_n(\omega)\in A\}
=
X_n^{-1}(A)\in\mathcal F.
\]

Значит \(X\) измеримо. Для конечномерного цилиндра:

\[
X^{-1}(\pi_{n_1,\ldots,n_k}^{-1}(B))
=
\{\omega:(X_{n_1}(\omega),\ldots,X_{n_k}(\omega))\in B\}\in\mathcal F,
\]

потому что вектор

\[
(X_{n_1},\ldots,X_{n_k})
\]

измерим.

Обратно, если задана случайная последовательность

\[
X:\Omega\to E^{\mathbb N_0},
\]

то отдельные координаты получаются как композиции:

\[
X_n=\pi_n\circ X.
\]

Так как \(\pi_n\) измерима, а \(X\) измеримо, то каждое \(X_n\) -
случайная величина. Значит одно случайное отображение в пространство
последовательностей задает семейство случайных величин.

Эта эквивалентность - центральная мысль билета:

\[
\{X_n\}_{n\ge0}
\quad\Longleftrightarrow\quad
X=(X_0,X_1,\ldots):\Omega\to E^{\mathbb N_0}.
\]

\subsubsection{2.6. Закон процесса и конечномерные
распределения}\label{ux437ux430ux43aux43eux43d-ux43fux440ux43eux446ux435ux441ux441ux430-ux438-ux43aux43eux43dux435ux447ux43dux43eux43cux435ux440ux43dux44bux435-ux440ux430ux441ux43fux440ux435ux434ux435ux43bux435ux43dux438ux44f}

Теперь закон процесса. Как и для любого случайного элемента из Билета
06, распределение случайной последовательности - это образ вероятностной
меры:

\[
P_X=P\circ X^{-1}.
\]

То есть для \(A\in\mathcal E^{\mathbb N_0}\):

\[
P_X(A)=P\{\omega:X(\omega)\in A\}.
\]

Это мера на пространстве последовательностей. Именно она называется
законом процесса.

Конечномерные распределения получаются из закона процесса проекциями.
Для \(n_1,\ldots,n_k\):

\[
P_{n_1,\ldots,n_k}(B)
=
P\{(X_{n_1},\ldots,X_{n_k})\in B\}.
\]

Через закон \(P_X\) это записывается так:

\[
P_{n_1,\ldots,n_k}(B)
=
P_X(\pi_{n_1,\ldots,n_k}^{-1}(B)).
\]

Иными словами, конечномерные распределения - это маргиналы меры \(P_X\).

\subsubsection{2.7. Связь с теоремой Колмогорова и канонический
процесс}\label{ux441ux432ux44fux437ux44c-ux441-ux442ux435ux43eux440ux435ux43cux43eux439-ux43aux43eux43bux43cux43eux433ux43eux440ux43eux432ux430-ux438-ux43aux430ux43dux43eux43dux438ux447ux435ux441ux43aux438ux439-ux43fux440ux43eux446ux435ux441ux441}

Наоборот, если заданы все конечномерные распределения и они согласованы,
то по Билету 15 существует мера \(P_X\) на \(E^{\mathbb N_0}\), для
которой эти распределения являются маргиналами. Тогда на каноническом
пространстве можно взять координатный процесс

\[
X_n(x)=x_n,
\qquad
x=(x_0,x_1,\ldots).
\]

Так получается случайная последовательность с заданными конечномерными
распределениями.

\subsubsection{2.8. Различие между состоянием, траекторией и
законом}\label{ux440ux430ux437ux43bux438ux447ux438ux435-ux43cux435ux436ux434ux443-ux441ux43eux441ux442ux43eux44fux43dux438ux435ux43c-ux442ux440ux430ux435ux43aux442ux43eux440ux438ux435ux439-ux438-ux437ux430ux43aux43eux43dux43eux43c}

Очень важно различать три объекта:

\begin{enumerate}
\def\labelenumi{\arabic{enumi}.}
\tightlist
\item
  \(X_n\) - состояние процесса в момент времени \(n\) как случайная
  величина;
\item
  \(X(\omega)\) - траектория, то есть последовательность значений при
  фиксированном исходе \(\omega\);
\item
  \(P_X\) - закон процесса, то есть вероятность на пространстве всех
  последовательностей.
\end{enumerate}

Траектория - не распределение. Распределение - мера на множестве
возможных траекторий.

\subsubsection{2.9. Определение закона через
цилиндры}\label{ux43eux43fux440ux435ux434ux435ux43bux435ux43dux438ux435-ux437ux430ux43aux43eux43dux430-ux447ux435ux440ux435ux437-ux446ux438ux43bux438ux43dux434ux440ux44b}

Для дискретного времени цилиндры особенно удобны. Например, если
\(E=\mathbb R\), то событие

\[
\{X_0\le a_0,\ X_1\le a_1,\ldots,X_n\le a_n\}
\]

соответствует цилиндру

\[
(-\infty,a_0]\times\cdots\times(-\infty,a_n]\times E\times E\times\cdots.
\]

Вероятность такого цилиндра равна конечномерной функции распределения:

\[
F_{0,\ldots,n}(a_0,\ldots,a_n)
=
P\{X_0\le a_0,\ldots,X_n\le a_n\}.
\]

Поэтому закон случайной последовательности полностью определяется
вероятностями цилиндров, а значит всеми конечномерными распределениями.

\subsection{3. Основные
определения}\label{ux43eux441ux43dux43eux432ux43dux44bux435-ux43eux43fux440ux435ux434ux435ux43bux435ux43dux438ux44f}

\subsubsection{Дискретный случайный
процесс}\label{ux434ux438ux441ux43aux440ux435ux442ux43dux44bux439-ux441ux43bux443ux447ux430ux439ux43dux44bux439-ux43fux440ux43eux446ux435ux441ux441}

Семейство случайных элементов

\[
\{X_n:n\in\mathbb N_0\}
\]

на одном вероятностном пространстве, где каждое

\[
X_n:\Omega\to E
\]

измеримо.

\subsubsection{Случайная
последовательность}\label{ux441ux43bux443ux447ux430ux439ux43dux430ux44f-ux43fux43eux441ux43bux435ux434ux43eux432ux430ux442ux435ux43bux44cux43dux43eux441ux442ux44c}

Случайный элемент со значениями в пространстве последовательностей:

\[
X:\Omega\to E^{\mathbb N_0},
\qquad
X(\omega)=(X_0(\omega),X_1(\omega),\ldots).
\]

\subsubsection{Пространство
последовательностей}\label{ux43fux440ux43eux441ux442ux440ux430ux43dux441ux442ux432ux43e-ux43fux43eux441ux43bux435ux434ux43eux432ux430ux442ux435ux43bux44cux43dux43eux441ux442ux435ux439}

\[
E^{\mathbb N_0}
=
\{(x_0,x_1,\ldots):x_n\in E\}.
\]

\subsubsection{\texorpdfstring{Произведенная
\(\sigma\)-алгебра}{Произведенная \textbackslash sigma-алгебра}}\label{ux43fux440ux43eux438ux437ux432ux435ux434ux435ux43dux43dux430ux44f-sigma-ux430ux43bux433ux435ux431ux440ux430}

\[
\mathcal E^{\mathbb N_0}
=
\sigma\{\pi_n^{-1}(A):n\ge0,\ A\in\mathcal E\}.
\]

\subsubsection{Координатная
проекция}\label{ux43aux43eux43eux440ux434ux438ux43dux430ux442ux43dux430ux44f-ux43fux440ux43eux435ux43aux446ux438ux44f}

\[
\pi_n(x)=x_n.
\]

\subsubsection{Цилиндрическое
событие}\label{ux446ux438ux43bux438ux43dux434ux440ux438ux447ux435ux441ux43aux43eux435-ux441ux43eux431ux44bux442ux438ux435}

\[
\pi_{n_1,\ldots,n_k}^{-1}(B),
\qquad
B\in\mathcal E^k.
\]

Это событие, зависящее от конечного числа координат.

\subsubsection{Траектория
процесса}\label{ux442ux440ux430ux435ux43aux442ux43eux440ux438ux44f-ux43fux440ux43eux446ux435ux441ux441ux430}

Для фиксированного исхода \(\omega\) последовательность

\[
X(\omega)=(X_0(\omega),X_1(\omega),\ldots).
\]

\subsubsection{Закон
процесса}\label{ux437ux430ux43aux43eux43d-ux43fux440ux43eux446ux435ux441ux441ux430}

\[
P_X=P\circ X^{-1}
\]

\begin{itemize}
\tightlist
\item
  распределение случайного элемента \(X\) на пространстве
  последовательностей.
\end{itemize}

\subsubsection{Конечномерное
распределение}\label{ux43aux43eux43dux435ux447ux43dux43eux43cux435ux440ux43dux43eux435-ux440ux430ux441ux43fux440ux435ux434ux435ux43bux435ux43dux438ux435}

\[
P_{n_1,\ldots,n_k}(B)
=
P\{(X_{n_1},\ldots,X_{n_k})\in B\}.
\]

\subsection{4. Основные теоремы и
свойства}\label{ux43eux441ux43dux43eux432ux43dux44bux435-ux442ux435ux43eux440ux435ux43cux44b-ux438-ux441ux432ux43eux439ux441ux442ux432ux430}

\subsubsection{Эквивалентность двух
описаний}\label{ux44dux43aux432ux438ux432ux430ux43bux435ux43dux442ux43dux43eux441ux442ux44c-ux434ux432ux443ux445-ux43eux43fux438ux441ux430ux43dux438ux439}

Последовательность случайных величин

\[
X_0,X_1,\ldots
\]

эквивалентна одному случайному отображению

\[
X:\Omega\to E^{\mathbb N_0},
\qquad
X(\omega)=(X_0(\omega),X_1(\omega),\ldots).
\]

\subsubsection{Критерий измеримости в пространство
последовательностей}\label{ux43aux440ux438ux442ux435ux440ux438ux439-ux438ux437ux43cux435ux440ux438ux43cux43eux441ux442ux438-ux432-ux43fux440ux43eux441ux442ux440ux430ux43dux441ux442ux432ux43e-ux43fux43eux441ux43bux435ux434ux43eux432ux430ux442ux435ux43bux44cux43dux43eux441ux442ux435ux439}

Отображение

\[
X:\Omega\to E^{\mathbb N_0}
\]

измеримо относительно произведенной \(\sigma\)-алгебры тогда и только
тогда, когда все координаты

\[
X_n=\pi_n\circ X
\]

измеримы.

\subsubsection{Закон процесса задается на пространстве
последовательностей}\label{ux437ux430ux43aux43eux43d-ux43fux440ux43eux446ux435ux441ux441ux430-ux437ux430ux434ux430ux435ux442ux441ux44f-ux43dux430-ux43fux440ux43eux441ux442ux440ux430ux43dux441ux442ux432ux435-ux43fux43eux441ux43bux435ux434ux43eux432ux430ux442ux435ux43bux44cux43dux43eux441ux442ux435ux439}

\[
P_X(A)=P\{X\in A\},
\qquad
A\in\mathcal E^{\mathbb N_0}.
\]

\subsubsection{Конечномерные распределения являются маргиналами
закона}\label{ux43aux43eux43dux435ux447ux43dux43eux43cux435ux440ux43dux44bux435-ux440ux430ux441ux43fux440ux435ux434ux435ux43bux435ux43dux438ux44f-ux44fux432ux43bux44fux44eux442ux441ux44f-ux43cux430ux440ux433ux438ux43dux430ux43bux430ux43cux438-ux437ux430ux43aux43eux43dux430}

\[
P_{n_1,\ldots,n_k}(B)
=
P_X(\pi_{n_1,\ldots,n_k}^{-1}(B)).
\]

\subsubsection{Закон процесса определяется
цилиндрами}\label{ux437ux430ux43aux43eux43d-ux43fux440ux43eux446ux435ux441ux441ux430-ux43eux43fux440ux435ux434ux435ux43bux44fux435ux442ux441ux44f-ux446ux438ux43bux438ux43dux434ux440ux430ux43cux438}

Если две вероятностные меры на \(E^{\mathbb N_0}\) совпадают на всех
цилиндрах, то они совпадают на всей произведенной \(\sigma\)-алгебре.

\subsubsection{Согласованные конечномерные распределения задают закон
последовательности}\label{ux441ux43eux433ux43bux430ux441ux43eux432ux430ux43dux43dux44bux435-ux43aux43eux43dux435ux447ux43dux43eux43cux435ux440ux43dux44bux435-ux440ux430ux441ux43fux440ux435ux434ux435ux43bux435ux43dux438ux44f-ux437ux430ux434ux430ux44eux442-ux437ux430ux43aux43eux43d-ux43fux43eux441ux43bux435ux434ux43eux432ux430ux442ux435ux43bux44cux43dux43eux441ux442ux438}

По теореме Колмогорова из Билета 15 согласованное семейство
конечномерных распределений задает меру на \(E^{\mathbb N_0}\), то есть
закон случайной последовательности.

\subsection{5. Примеры}\label{ux43fux440ux438ux43cux435ux440ux44b}

\subsubsection{Независимые броски
монеты}\label{ux43dux435ux437ux430ux432ux438ux441ux438ux43cux44bux435-ux431ux440ux43eux441ux43aux438-ux43cux43eux43dux435ux442ux44b}

Пусть

\[
E=\{0,1\},
\]

где \(1\) означает успех, а \(0\) - неуспех. Последовательность

\[
X_0,X_1,\ldots
\]

независимых бернуллиевских величин с параметром \(p\) имеет
конечномерные вероятности

\[
P\{X_0=x_0,\ldots,X_n=x_n\}
=
\prod_{k=0}^n p^{x_k}(1-p)^{1-x_k},
\qquad
x_k\in\{0,1\}.
\]

Ее закон - мера на пространстве

\[
\{0,1\}^{\mathbb N_0}.
\]

Траектория - бесконечная последовательность нулей и единиц.

\subsubsection{Случайное
блуждание}\label{ux441ux43bux443ux447ux430ux439ux43dux43eux435-ux431ux43bux443ux436ux434ux430ux43dux438ux435}

Пусть \(\xi_1,\xi_2,\ldots\) - независимые одинаково распределенные
приращения, например

\[
P\{\xi_k=1\}=P\{\xi_k=-1\}=\frac12.
\]

Определим

\[
S_0=0,
\qquad
S_n=\sum_{k=1}^n \xi_k.
\]

Тогда

\[
S=(S_0,S_1,S_2,\ldots)
\]

\begin{itemize}
\tightlist
\item
  случайная последовательность со значениями в
  \(\mathbb Z^{\mathbb N_0}\). Каждая траектория - путь блуждания.
\end{itemize}

\subsubsection{Марковская
цепь}\label{ux43cux430ux440ux43aux43eux432ux441ux43aux430ux44f-ux446ux435ux43fux44c}

Пусть \(E\) - счетное множество состояний, \(\pi\) - начальное
распределение, \(P=(p_{ij})\) - переходная матрица. Тогда конечномерные
вероятности имеют вид

\[
P\{X_0=i_0,\ldots,X_n=i_n\}
=
\pi_{i_0}p_{i_0i_1}\cdots p_{i_{n-1}i_n}.
\]

Это задает закон на \(E^{\mathbb N_0}\). Координатная последовательность
с этим законом является марковской цепью. Это связь с Билетом 22 и
Билетом 23.

\subsubsection{Стационарная
последовательность}\label{ux441ux442ux430ux446ux438ux43eux43dux430ux440ux43dux430ux44f-ux43fux43eux441ux43bux435ux434ux43eux432ux430ux442ux435ux43bux44cux43dux43eux441ux442ux44c}

Случайная последовательность называется стационарной в узком смысле,
если ее конечномерные распределения не меняются при сдвиге времени:

\[
(X_{n_1+h},\ldots,X_{n_k+h})
\stackrel d=
(X_{n_1},\ldots,X_{n_k})
\]

для всех допустимых \(h\). В терминах закона на \(E^{\mathbb N_0}\) это
означает инвариантность закона относительно сдвига последовательности.
Это связь с Билетом 17.

\subsubsection{Пример различия траектории и
закона}\label{ux43fux440ux438ux43cux435ux440-ux440ux430ux437ux43bux438ux447ux438ux44f-ux442ux440ux430ux435ux43aux442ux43eux440ux438ux438-ux438-ux437ux430ux43aux43eux43dux430}

Если бросать монету бесконечно много раз, то конкретная траектория может
быть

\[
(1,0,0,1,1,0,\ldots).
\]

Но закон процесса - не эта последовательность, а мера на множестве всех
бесконечных последовательностей из нулей и единиц. Эта мера говорит, с
какими вероятностями возникают цилиндры, например событие

\[
\{x_0=1,\ x_1=0,\ x_2=0\}.
\]

\subsection{6. Схемы доказательств /
обоснований}\label{ux441ux445ux435ux43cux44b-ux434ux43eux43aux430ux437ux430ux442ux435ux43bux44cux441ux442ux432-ux43eux431ux43eux441ux43dux43eux432ux430ux43dux438ux439}

\subsubsection{Из последовательности случайных величин получается
случайная
последовательность}\label{ux438ux437-ux43fux43eux441ux43bux435ux434ux43eux432ux430ux442ux435ux43bux44cux43dux43eux441ux442ux438-ux441ux43bux443ux447ux430ux439ux43dux44bux445-ux432ux435ux43bux438ux447ux438ux43d-ux43fux43eux43bux443ux447ux430ux435ux442ux441ux44f-ux441ux43bux443ux447ux430ux439ux43dux430ux44f-ux43fux43eux441ux43bux435ux434ux43eux432ux430ux442ux435ux43bux44cux43dux43eux441ux442ux44c}

Пусть все \(X_n\) измеримы. Определим

\[
X(\omega)=(X_0(\omega),X_1(\omega),\ldots).
\]

Чтобы доказать измеримость \(X\), достаточно проверить прообразы
образующих цилиндров:

\[
X^{-1}(\pi_n^{-1}(A))
=
\{\omega:X_n(\omega)\in A\}
\in\mathcal F.
\]

Значит \(X\) измеримо как отображение в \(E^{\mathbb N_0}\).

\subsubsection{Из случайной последовательности получаются
координаты}\label{ux438ux437-ux441ux43bux443ux447ux430ux439ux43dux43eux439-ux43fux43eux441ux43bux435ux434ux43eux432ux430ux442ux435ux43bux44cux43dux43eux441ux442ux438-ux43fux43eux43bux443ux447ux430ux44eux442ux441ux44f-ux43aux43eux43eux440ux434ux438ux43dux430ux442ux44b}

Пусть \(X:\Omega\to E^{\mathbb N_0}\) измеримо. Тогда

\[
X_n=\pi_n\circ X.
\]

Поскольку координатные проекции \(\pi_n\) измеримы, каждая \(X_n\)
измерима. Значит координаты случайной последовательности являются
случайными величинами.

\subsubsection{Почему закон определяется конечномерными
распределениями}\label{ux43fux43eux447ux435ux43cux443-ux437ux430ux43aux43eux43d-ux43eux43fux440ux435ux434ux435ux43bux44fux435ux442ux441ux44f-ux43aux43eux43dux435ux447ux43dux43eux43cux435ux440ux43dux44bux43cux438-ux440ux430ux441ux43fux440ux435ux434ux435ux43bux435ux43dux438ux44fux43cux438}

Конечномерные распределения задают вероятности всех цилиндров:

\[
P_X(\pi_{n_1,\ldots,n_k}^{-1}(B))
=
P\{(X_{n_1},\ldots,X_{n_k})\in B\}.
\]

Цилиндры порождают произведенную \(\sigma\)-алгебру. Поэтому мера на
пространстве последовательностей определяется своими значениями на
цилиндрах.

\subsubsection{Как используется теорема
Колмогорова}\label{ux43aux430ux43a-ux438ux441ux43fux43eux43bux44cux437ux443ux435ux442ux441ux44f-ux442ux435ux43eux440ux435ux43cux430-ux43aux43eux43bux43cux43eux433ux43eux440ux43eux432ux430}

Если процесс еще не задан, но есть согласованные конечномерные
распределения, то теорема Колмогорова строит меру \(P\) на
\(E^{\mathbb N_0}\). После этого можно взять канонический процесс

\[
X_n(x)=x_n.
\]

Его закон равен \(P\), а конечномерные распределения совпадают с
заданными.

\subsection{7. Что написать на
доске}\label{ux447ux442ux43e-ux43dux430ux43fux438ux441ux430ux442ux44c-ux43dux430-ux434ux43eux441ux43aux435}

\begin{enumerate}
\def\labelenumi{\arabic{enumi}.}
\tightlist
\item
  Две точки зрения:
\end{enumerate}

\[
\{X_n\}_{n\ge0}
\quad\Longleftrightarrow\quad
X:\Omega\to E^{\mathbb N_0}.
\]

\begin{enumerate}
\def\labelenumi{\arabic{enumi}.}
\setcounter{enumi}{1}
\tightlist
\item
  Случайная последовательность:
\end{enumerate}

\[
X(\omega)=(X_0(\omega),X_1(\omega),\ldots).
\]

\begin{enumerate}
\def\labelenumi{\arabic{enumi}.}
\setcounter{enumi}{2}
\tightlist
\item
  Пространство последовательностей:
\end{enumerate}

\[
E^{\mathbb N_0}.
\]

\begin{enumerate}
\def\labelenumi{\arabic{enumi}.}
\setcounter{enumi}{3}
\tightlist
\item
  Координатные проекции:
\end{enumerate}

\[
\pi_n(x)=x_n,
\qquad
X_n=\pi_n\circ X.
\]

\begin{enumerate}
\def\labelenumi{\arabic{enumi}.}
\setcounter{enumi}{4}
\tightlist
\item
  Произведенная \(\sigma\)-алгебра:
\end{enumerate}

\[
\mathcal E^{\mathbb N_0}
=
\sigma\{\pi_n^{-1}(A):A\in\mathcal E,\ n\ge0\}.
\]

\begin{enumerate}
\def\labelenumi{\arabic{enumi}.}
\setcounter{enumi}{5}
\tightlist
\item
  Закон процесса:
\end{enumerate}

\[
P_X=P\circ X^{-1}.
\]

\begin{enumerate}
\def\labelenumi{\arabic{enumi}.}
\setcounter{enumi}{6}
\tightlist
\item
  Конечномерные распределения как маргиналы:
\end{enumerate}

\[
P_{n_1,\ldots,n_k}(B)
=
P_X(\pi_{n_1,\ldots,n_k}^{-1}(B)).
\]

\subsection{8. Типичные дополнительные
вопросы}\label{ux442ux438ux43fux438ux447ux43dux44bux435-ux434ux43eux43fux43eux43bux43dux438ux442ux435ux43bux44cux43dux44bux435-ux432ux43eux43fux440ux43eux441ux44b}

\textbf{Что является главным объектом билета?}

Дискретный случайный процесс, который можно рассматривать либо как
последовательность случайных величин \(X_0,X_1,\ldots\), либо как один
случайный элемент \(X\) со значениями в пространстве последовательностей
\(E^{\mathbb N_0}\).

\textbf{В чем разница между \(X_n\) и \(X(\omega)\)?}

\(X_n\) - случайная величина, отвечающая моменту времени \(n\). А
\(X(\omega)\) - вся траектория при фиксированном исходе \(\omega\):

\[
X(\omega)=(X_0(\omega),X_1(\omega),\ldots).
\]

\textbf{Что такое пространство последовательностей?}

Это множество всех возможных траекторий:

\[
E^{\mathbb N_0}=\{(x_0,x_1,\ldots):x_n\in E\}.
\]

На нем используется произведенная \(\sigma\)-алгебра, порожденная
координатными цилиндрами.

\textbf{Что такое цилиндрическое событие?}

Это событие, которое задает условия только на конечное число координат:

\[
\{x:(x_{n_1},\ldots,x_{n_k})\in B\}.
\]

Например, \(\{x_0\in A_0,\ x_3\in A_3\}\).

\textbf{Почему отображение \(X:\Omega\to E^{\mathbb N_0}\) измеримо,
если все координаты \(X_n\) измеримы?}

Потому что произведенная \(\sigma\)-алгебра порождена цилиндрами, а
прообраз координатного цилиндра равен

\[
X^{-1}(\pi_n^{-1}(A))=X_n^{-1}(A)\in\mathcal F.
\]

\textbf{Что такое закон процесса?}

Это распределение случайной последовательности:

\[
P_X(A)=P\{X\in A\},
\qquad
A\in\mathcal E^{\mathbb N_0}.
\]

Это мера на множестве всех возможных траекторий.

\textbf{Как конечномерные распределения связаны с законом процесса?}

Они являются маргиналами закона:

\[
P_{n_1,\ldots,n_k}(B)
=
P_X(\pi_{n_1,\ldots,n_k}^{-1}(B)).
\]

\textbf{Почему нельзя путать траекторию и распределение?}

Траектория - это один элемент \(E^{\mathbb N_0}\), то есть конкретная
последовательность значений. Распределение - это вероятностная мера на
множестве всех таких последовательностей.

\textbf{Как билет связан с Билетом 15?}

В Билете 15 доказывается, что согласованные конечномерные распределения
задают меру на пространстве траекторий. В этом билете для дискретного
времени это пространство имеет вид \(E^{\mathbb N_0}\).

\textbf{Как билет связан с марковскими цепями?}

Марковская цепь - это частный случай случайной последовательности. Ее
закон на \(E^{\mathbb N_0}\) задается начальным распределением и
переходными вероятностями:

\[
P\{X_0=i_0,\ldots,X_n=i_n\}
=
\pi_{i_0}p_{i_0i_1}\cdots p_{i_{n-1}i_n}.
\]

\subsection{9. Типичные
ошибки}\label{ux442ux438ux43fux438ux447ux43dux44bux435-ux43eux448ux438ux431ux43aux438}

\begin{itemize}
\tightlist
\item
  Путать время \(n\) и исход \(\omega\).
\item
  Путать случайную величину \(X_n\) и значение \(X_n(\omega)\).
\item
  Путать траекторию \(X(\omega)\) и закон процесса \(P_X\).
\item
  Забывать, что все \(X_n\) должны быть заданы на одном вероятностном
  пространстве.
\item
  Забывать произведенную \(\sigma\)-алгебру на \(E^{\mathbb N_0}\).
\item
  Считать, что одномерные распределения \(P_{X_n}\) задают закон
  процесса. Нужны все конечномерные распределения или мера на
  пространстве последовательностей.
\item
  Не различать семейство случайных величин и канонический координатный
  процесс на \(E^{\mathbb N_0}\).
\end{itemize}

\subsection{10. Связи с другими
билетами}\label{ux441ux432ux44fux437ux438-ux441-ux434ux440ux443ux433ux438ux43cux438-ux431ux438ux43bux435ux442ux430ux43cux438}

\begin{itemize}
\tightlist
\item
  Билет 05: случайные элементы и измеримые отображения.
\item
  Билет 06: распределение случайного элемента и конечномерные
  распределения.
\item
  Билет 15: согласованные конечномерные распределения задают меру на
  пространстве траекторий.
\item
  Билет 17: стационарность и эргодичность формулируются для случайных
  последовательностей.
\item
  Билет 19: случайные блуждания - примеры случайных последовательностей
  с независимыми приращениями.
\item
  Билет 22 и Билет 23: марковские цепи - случайные последовательности со
  специальной зависимостью.
\item
  Билет 24 и Билет 25: корреляционная функция и спектральные
  характеристики относятся к последовательностям второго порядка.
\end{itemize}

\subsection{11. Минимум на
удовлетворительно}\label{ux43cux438ux43dux438ux43cux443ux43c-ux43dux430-ux443ux434ux43eux432ux43bux435ux442ux432ux43eux440ux438ux442ux435ux43bux44cux43dux43e}

Нужно сказать:

\begin{enumerate}
\def\labelenumi{\arabic{enumi}.}
\tightlist
\item
  Дискретный процесс - это последовательность случайных величин:
\end{enumerate}

\[
X_0,X_1,\ldots
\]

\begin{enumerate}
\def\labelenumi{\arabic{enumi}.}
\setcounter{enumi}{1}
\tightlist
\item
  Его можно рассматривать как одну случайную последовательность:
\end{enumerate}

\[
X:\Omega\to E^{\mathbb N_0},
\qquad
X(\omega)=(X_0(\omega),X_1(\omega),\ldots).
\]

\begin{enumerate}
\def\labelenumi{\arabic{enumi}.}
\setcounter{enumi}{2}
\tightlist
\item
  Закон процесса:
\end{enumerate}

\[
P_X=P\circ X^{-1}
\]

на пространстве \(E^{\mathbb N_0}\).

\begin{enumerate}
\def\labelenumi{\arabic{enumi}.}
\setcounter{enumi}{3}
\item
  Конечномерные распределения - маргиналы этого закона.
\item
  Главная ловушка: траектория - не распределение.
\end{enumerate}

\subsection{12. Минимум на
хорошо}\label{ux43cux438ux43dux438ux43cux443ux43c-ux43dux430-ux445ux43eux440ux43eux448ux43e}

Помимо минимума, нужно объяснить:

\begin{itemize}
\tightlist
\item
  что такое произведенная \(\sigma\)-алгебра;
\item
  что такое цилиндрическое событие;
\item
  почему измеримость отображения в \(E^{\mathbb N_0}\) проверяется по
  координатам;
\item
  как из закона процесса получить конечномерные распределения;
\item
  как по согласованным конечномерным распределениям построить закон
  последовательности через теорему Колмогорова;
\item
  привести пример: iid-последовательность, случайное блуждание или
  марковская цепь.
\end{itemize}

\subsection{13. Что добавить для
отлично}\label{ux447ux442ux43e-ux434ux43eux431ux430ux432ux438ux442ux44c-ux434ux43bux44f-ux43eux442ux43bux438ux447ux43dux43e}

Для сильного ответа стоит добавить:

\begin{itemize}
\tightlist
\item
  четкое различие между \(X_n\), \(X_n(\omega)\), \(X(\omega)\) и
  \(P_X\);
\item
  формулу для цилиндров и конечномерных проекций;
\item
  доказательство эквивалентности двух точек зрения через измеримость
  координат;
\item
  объяснение, почему закон на \(E^{\mathbb N_0}\) определяется
  цилиндрами;
\item
  пример марковской цепи через \(\pi\) и \(P\);
\item
  связь со стационарностью как инвариантностью закона относительно
  сдвига.
\end{itemize}

\subsection{14. Устная версия ответа на 5
минут}\label{ux443ux441ux442ux43dux430ux44f-ux432ux435ux440ux441ux438ux44f-ux43eux442ux432ux435ux442ux430-ux43dux430-5-ux43cux438ux43dux443ux442}

\begin{enumerate}
\def\labelenumi{\arabic{enumi}.}
\tightlist
\item
  Дискретный случайный процесс - это семейство случайных величин
  \(X_0,X_1,\ldots\) на одном вероятностном пространстве.
\item
  Его можно рассматривать как одно отображение
\end{enumerate}

\[
X:\Omega\to E^{\mathbb N_0},
\qquad
X(\omega)=(X_0(\omega),X_1(\omega),\ldots).
\]

\begin{enumerate}
\def\labelenumi{\arabic{enumi}.}
\setcounter{enumi}{2}
\tightlist
\item
  Пространство \(E^{\mathbb N_0}\) - это пространство всех траекторий.
\item
  На нем берется произведенная \(\sigma\)-алгебра, порожденная
  координатными проекциями
\end{enumerate}

\[
\pi_n(x)=x_n.
\]

\begin{enumerate}
\def\labelenumi{\arabic{enumi}.}
\setcounter{enumi}{4}
\tightlist
\item
  Цилиндры - события, зависящие от конечного числа координат:
\end{enumerate}

\[
\{x:(x_{n_1},\ldots,x_{n_k})\in B\}.
\]

\begin{enumerate}
\def\labelenumi{\arabic{enumi}.}
\setcounter{enumi}{5}
\tightlist
\item
  Если все \(X_n\) измеримы, то отображение \(X=(X_0,X_1,\ldots)\)
  измеримо, потому что прообразы цилиндров измеримы.
\item
  Обратно, если \(X\) - случайный элемент пространства
  последовательностей, то его координаты
\end{enumerate}

\[
X_n=\pi_n\circ X
\]

являются случайными величинами. 8. Закон процесса - это распределение
\(X\):

\[
P_X=P\circ X^{-1}.
\]

\begin{enumerate}
\def\labelenumi{\arabic{enumi}.}
\setcounter{enumi}{8}
\tightlist
\item
  Конечномерные распределения - маргиналы этого закона:
\end{enumerate}

\[
P_{n_1,\ldots,n_k}(B)
=
P_X(\pi_{n_1,\ldots,n_k}^{-1}(B)).
\]

\begin{enumerate}
\def\labelenumi{\arabic{enumi}.}
\setcounter{enumi}{9}
\tightlist
\item
  По теореме Колмогорова согласованные конечномерные распределения
  задают такой закон на \(E^{\mathbb N_0}\).
\item
  Примеры: независимые броски монеты, случайное блуждание, марковская
  цепь.
\item
  Главная ошибка - путать конкретную траекторию с распределением на
  пространстве траекторий.
\end{enumerate}

\subsection{15. Если осталось 2 минуты перед
экзаменом}\label{ux435ux441ux43bux438-ux43eux441ux442ux430ux43bux43eux441ux44c-2-ux43cux438ux43dux443ux442ux44b-ux43fux435ux440ux435ux434-ux44dux43aux437ux430ux43cux435ux43dux43eux43c}

Запомнить главную эквивалентность:

\[
\{X_n\}_{n\ge0}
\quad\Longleftrightarrow\quad
X:\Omega\to E^{\mathbb N_0},
\qquad
X(\omega)=(X_0(\omega),X_1(\omega),\ldots).
\]

Координаты:

\[
X_n=\pi_n\circ X,
\qquad
\pi_n(x)=x_n.
\]

Закон процесса:

\[
P_X=P\circ X^{-1}.
\]

Конечномерные распределения:

\[
P_{n_1,\ldots,n_k}(B)
=
P_X(\pi_{n_1,\ldots,n_k}^{-1}(B)).
\]

Главная ловушка: \(X(\omega)\) - это одна траектория, а \(P_X\) - мера
на множестве всех траекторий.

\newpage

\end{document}
